(1) 定态方程 \(-\hbar^2\psi''/(2m)-g\delta(x)\psi=E\psi\)。在原点两侧取无穷小区间积分，得到 \(\psi'(0^+)-\psi'(0^-)=-2mg\psi(0)/\hbar^2\)，而波函数本身必须连续，否则二阶导会产生不允许的 \(\delta'\) 分布。
(2) 对 \(E<0\)，阱外衰减解为 \(\psi(x)=A e^{-\kappa|x|}\)，\(\kappa=\sqrt{-2mE}/\hbar\)。导数跳跃给 \(-2\kappa A=-2mgA/\hbar^2\)，所以 \(\kappa=mg/\hbar^2\)，\(E_b=-mg^2/(2\hbar^2)\)；归一积分 \(\int e^{-2\kappa|x|}dx=1/\kappa\) 得 \(A=\sqrt\kappa\)，仅有一个束缚态。
(3) 正能量波数 \(k=\sqrt{2mE}/\hbar\)，设 \(\psi(x<0)=e^{ikx}+r e^{-ikx}\)、\(\psi(x>0)=t e^{ikx}\)。连续给 \(t=1+r\)；导数跳跃为 \(2ik r=-2\kappa t\)，解得 \(r=-\kappa/(\kappa+ik)\)、\(t=ik/(\kappa+ik)\)，\(\kappa=mg/\hbar^2\)。
(4) 介质两侧波数相同，概率流比直接为 \(R=|r|^2=\kappa^2/(k^2+\kappa^2)\)、\(T=|t|^2=k^2/(k^2+\kappa^2)\)，故 \(R+T=1\)。\(g\to0\) 时 \(\kappa\to0,T\to1\)；固定吸引强度且 \(k\to0^+\) 时 \(R\to1\)，虽然势为吸引型，低能散射仍趋全反射。
